\documentclass[11pt]{article}
\usepackage[margin=1in]{geometry}
\usepackage{graphicx}
\usepackage{hyperref}
\usepackage{parskip}

\title{User Guide: Mores Creek Summit UAVSAR/Lidar Datacube\\
\large \texttt{ZARR/mores\_creek\_summit.zarr} and the modeling-team subset
\texttt{mcs\_top\_pairs\_subset.nc}}
\author{HP Marshall, Boise State University}
\date{September 2026}

\begin{document}
\maketitle

\section{Overview}

The datacube collects every SnowEx UAVSAR interferometric product and airborne lidar
product over Mores Creek Summit (MCS), Idaho, on one shared analysis grid, together with
the derived quantities needed to turn interferometric phase into snow water equivalent
change ($\Delta$SWE): local incidence angle, atmospheric phase delay, and the $\Delta$SWE
retrievals themselves.

\begin{itemize}
\item \textbf{Format:} Zarr v2, consolidated metadata (9.4\,GB on disk, 15\,GB logical).
\item \textbf{Grid:} 3\,m, UTM zone 11N (EPSG:32611); 2801 rows $\times$ 2624 columns;
      affine transform $(3, 0, 601558\,/\,0, -3, 4870872.5)$; derived from the native
      0.5\,m MCS lidar grid. Northing (\texttt{y}) is stored north~$\to$~south.
\item \textbf{Bounds (lon/lat):} $-115.7352$ to $-115.6355$, $43.9072$ to $43.9840$.
\item \textbf{Missing data:} \texttt{NaN} everywhere, including outside each product's
      footprint.
\end{itemize}

\section{Dimensions and coordinates}

\noindent\resizebox{\textwidth}{!}{\begin{tabular}{llll}
\hline
name & size / dims & dtype & meaning \\
\hline
flight (dim) & 20 & -- & unique flight index (line, date) \\
y (dim) & 2801 & -- & UTM northing of cell centers (m), stored north to south \\
x (dim) & 2624 & -- & UTM easting of cell centers (m, EPSG:32611) \\
pair (dim) & 22 & -- & InSAR pair index (order of uavsar\_pair\_manifest\_mcs.json) \\
lidar\_time (dim) & 13 & -- & MCS\_Lidar acquisition date (DTM/DSM/CHM) \\
pol (dim) & 4 & -- & polarization (HH, HV, VH, VV) \\
line (dim) & 2 & -- & UAVSAR flight line id (05208 NE-looking, 23205 SW-looking) \\
sd\_time (dim) & 15 & -- & snow-depth acquisition date \\
flight\_date & coord on flight & str(10) & acquisition date of each flight \\
flight\_line\_id & coord on flight & str(5) & flight line of each flight \\
heading & coord on pair & int64 & aircraft heading of each pair (deg) \\
lidar\_time & coord on lidar\_time & datetime & MCS\_Lidar acquisition date (DTM/DSM/CHM) \\
line & coord on line & str(5) & UAVSAR flight line id (05208 NE-looking, 23205 SW-looking) \\
line\_id & coord on pair & str(5) & flight line of each pair \\
pair\_date1 & coord on pair & datetime & pair reference (earlier) acquisition date \\
pair\_date2 & coord on pair & datetime & pair repeat (later) acquisition date \\
pol & coord on pol & str(2) & polarization (HH, HV, VH, VV) \\
sd\_source & coord on sd\_time & str(14) & snow-depth product (SNEX20\_QSI\_SD or SNEX\_MCS\_Lidar) \\
sd\_time & coord on sd\_time & datetime & snow-depth acquisition date \\
temporal\_baseline\_days & coord on pair & int32 & days between the pair's acquisitions \\
x & coord on x & float64 & UTM easting of cell centers (m, EPSG:32611) \\
y & coord on y & float64 & UTM northing of cell centers (m), stored north to south \\
\hline
\end{tabular}
}

Two time axes coexist deliberately: \texttt{sd\_time} (15 snow-depth dates, 2020--2025,
QSI + MCS\_Lidar) and \texttt{lidar\_time} (13 MCS\_Lidar dates, 2022--2025, the only
collection providing DTM/DSM/CHM). The 22 pairs and 20 flights cover winters 2019--20
(line 23205 only) and 2020--21 (both lines).

\section{Data variables}

\noindent\resizebox{\textwidth}{!}{\begin{tabular}{p{2.6cm}p{3.1cm}p{1.5cm}p{2.0cm}p{5.6cm}}
\hline
variable & dims & dtype & chunks & units / source \\
\hline
atm\_delay & (flight, y, x) & float64 & (3, 351, 328) & radians; ERA5 delay below aircraft altitude, per flight (compute\_atm\_delay.py) \\
atm\_delay\_diff & (pair, y, x) & float64 & (3, 351, 328) & radians; per-pair delay difference (date2$-$date1) (compute\_atm\_delay\_diff.py) \\
chm & (lidar\_time, y, x) & float32 & (1, 512, 512) & m; MCS\_Lidar canopy height (build\_mcs\_datacube.py) \\
coherence & (pair, pol, y, x) & float32 & (1, 1, 512, 512) & unitless [0,1]; ASF .cor.grd (build\_mcs\_datacube.py) \\
dsm & (lidar\_time, y, x) & float32 & (1, 512, 512) & m; MCS\_Lidar surface DSM (build\_mcs\_datacube.py) \\
dswe & (pair, pol, y, x) & float32 & (1, 1, 512, 512) & m w.e.; Guneriussen-2001 inversion, $\rho$=250, SNOTEL-anchored, pair 8 deramped (compute\_dswe.py) \\
dtm & (lidar\_time, y, x) & float32 & (1, 512, 512) & m; MCS\_Lidar bare-earth DTM (build\_mcs\_datacube.py) \\
lia & (line, y, x) & float64 & (1, 351, 328) & degrees; JPL .llh/.lkv look vectors + 2023-02-09 DTM (compute\_local\_incidence\_angle.py) \\
sd & (sd\_time, y, x) & float32 & (1, 512, 512) & m; QSI + MCS\_Lidar snow depth, 0.5\,m block-averaged to 3\,m (build\_mcs\_datacube.py) \\
unwrapped\_phase & (pair, pol, y, x) & float32 & (1, 1, 512, 512) & radians; ASF .unw.grd; pairs 9/12/14/20/21 SNAPHU-unwrapped in-house (build\_mcs\_datacube.py, unwrap\_missing\_pairs.py) \\
wrapped\_phase & (pair, pol, y, x) & complex64 & (1, 1, 512, 512) & complex (unitless); ASF INTERFEROMETRY\_GRD, complex-domain regrid (build\_mcs\_datacube.py) \\
\hline
\end{tabular}
}

Variable-level attributes record provenance: \texttt{unwrapped\_phase} carries
\texttt{custom\_unwrapped\_pairs = [9, 12, 14, 20, 21]} (pairs unwrapped in-house with
SNAPHU because ASF ships no \texttt{.unw.grd}) and \texttt{dswe} carries the density
assumption, mask rule, and SNOTEL anchoring description.

\section{Dataset counts}

Number of 2-D $(y, x)$ maps per variable:

\medskip
\noindent\begin{tabular}{lrl}
\hline
variable & 2-D maps & composition \\
\hline
atm\_delay & 20 & 20 flight \\
atm\_delay\_diff & 22 & 22 pair \\
chm & 13 & 13 lidar\_time \\
coherence & 88 & 22 pair $\times$ 4 pol \\
dsm & 13 & 13 lidar\_time \\
dswe & 88 & 22 pair $\times$ 4 pol \\
dtm & 13 & 13 lidar\_time \\
lia & 2 & 2 line \\
sd & 15 & 15 sd\_time \\
unwrapped\_phase & 88 & 22 pair $\times$ 4 pol \\
wrapped\_phase & 88 & 22 pair $\times$ 4 pol \\
\hline
\end{tabular}


\section{The \texttt{dswe} retrieval}

$\Delta$SWE (m water equivalent, positive = SWE gain from \texttt{pair\_date1} to
\texttt{pair\_date2}) is retrieved per pair and polarization as follows: atmospheric
correction $\varphi = \texttt{unwrapped\_phase} - \texttt{atm\_delay\_diff}$; for pair~8 a
planar ramp (fit jointly with an elevation term that is deliberately \emph{kept}) is
removed; inversion via the Guneriussen et al.\ (2001) relation with the per-line local
incidence angle and new-snow density 250\,kg\,m$^{-3}$; masking where coherence $< 0.35$
or $\cos(\mathrm{LIA}) \le 0.05$; and anchoring of each map's unknown constant so the
90\,m window at the Mores Creek Summit SNOTEL (637:ID:SNTL, 43.932$^\circ$N,
$-115.666^\circ$E) matches the station's measured interval $\Delta$SWE. Pairs 0, 1 and 13
could not be anchored (no coherent unwrapped data within 600\,m of the station) and
retain an arbitrary constant; the SNAPHU-unwrapped pairs may contain several connected
components with independent constants. Validation (closure, cross-polarization, SNOTEL,
lidar) is documented in \texttt{dswe\_validation.md} and summarized in the companion
report.

\section{Reading the store}

Python (conda env \texttt{myenv}; see \texttt{environment.yml}):
\begin{verbatim}
import xarray as xr
ds = xr.open_zarr("ZARR/mores_creek_summit.zarr")
dswe = ds["dswe"].sel(pol="HH").isel(pair=8).load()   # one map, lazily chunked
\end{verbatim}
Note: set \texttt{PROJ\_DATA=/opt/anaconda3/envs/myenv/share/proj} before any
reprojection work (see AGENTS.md). MATLAB has no Zarr reader on this system; use the
NetCDF mirror \texttt{ZARR/mcs\_gui.nc} (regenerated by \texttt{export\_mcs\_gui\_nc.py})
or the interactive viewer \texttt{mcs\_gui.m}.

\section{The modeling-team subset: \texttt{mcs\_top\_pairs\_subset.nc}}

A 1.2\,GB NetCDF4 file (regenerated by \texttt{export\_subset\_nc.py}) containing only
what the snow modeling team needs, on the same 3\,m grid:

\begin{itemize}
\item \textbf{pair} dimension of 6, in rank order --- master-Zarr pairs 8, 2, 1, 5, 6
      (ranks 1--5) and pair 4 (zero-change example); the \texttt{rank\_note} coordinate
      labels each. All six carry original ASF unwrapped phase (none needed custom
      unwrapping).
\item Per pair: \texttt{dswe} (4 pols), \texttt{coherence} (4 pols, for quality
      screening), \texttt{unwrapped\_phase} (4 pols) and \texttt{atm\_delay\_diff}
      (for transparency --- already subtracted inside \texttt{dswe}).
\item Site layers: \texttt{sd} for the two QSI lidar surveys (2020-02-09, 2021-03-15),
      \texttt{dem} (2023-02-09 lidar DTM), \texttt{lia} (both headings).
\item Dates and labels are plain strings; all provenance is in the global attributes.
\end{itemize}

MATLAB: \verb|ncread('mcs_top_pairs_subset.nc','dswe',[1 1 1 1],[Inf Inf 1 1])| reads the
first pair's HH map (MATLAB returns dimensions reversed: $x, y, \mathrm{pol},
\mathrm{pair}$). Python: \verb|xr.open_dataset(...)|.

\end{document}
